The next class of functions $T: N \rightarrow N$, beyond the polynomials and having nice closure properties is the class of exponential functions. The complexity
classes \textbf{EXP} and \textbf{FEXP} are defined as follows:
$$ \textbf{EXP} = \bigcup_{c \geq 1} \textbf{DTIME}(2^{n^c}) \textbf{ FEXP} = \bigcup_{c \geq 1} \textbf{FDTIME}(2^{n^c})$$
The tasks in these classes can be solved mechanically, but possibly cannot be solved efficiently. Of course, the previous complexity classes \textbf{P} and 
\textbf{FP}, respectively for languages and functions, are subsets of the ones introduced.
$$ \textbf{P} \subseteq \textbf{EXP} \textbf{ FP} \subseteq \textbf{FEXP}$$
\begin{definition}[title = Theorem]
    The two inclusions above are strict.
\end{definition}
The theorem seems to be quite weird, but what it's saying is quite easy to understand:
\begin{itemize}
    \renewcommand{\labelitemi}{-}
    \item Some problems are in \textbf{EXP} complexity class but not in \textbf{P}.
    \item Some functions are in \textbf{FEXP} complexity class but not in \textbf{FP}.
\end{itemize} 